1. Change of metric signature
We now take:
\[
\eta_{\mu\nu} = \mathrm{diag}(+1, -1, -1, -1)
\]
So:
\[
u^2 = \eta_{\mu\nu} u^\mu u^\nu = (u^0)^2 - \mathbf{u}^2
\]
Timelike vectors satisfy:
\[
u^2 > 0
\]
2. Relativistic particle action
With this signature, the natural invariant length is:
\[
ds = \sqrt{u^2}\, d\lambda
\]
So the action becomes:
\[
S = -m \int d\lambda \sqrt{u^2}
\]
(Some authors omit the minus sign — we’ll address that.)
3. Compute the canonical momentum
Start from:
\[
L = -m \sqrt{u^2}
\]
Then:
\[
p_\mu = \frac{\partial L}{\partial u^\mu}
\]
Compute explicitly:
\[
p_\mu
= -m \frac{1}{2} (u^2)^{-1/2} (2 \eta_{\mu\nu} u^\nu)
= -m \frac{u_\mu}{\sqrt{u^2}}
\]
So again:
\[
p_\mu = -m \frac{u_\mu}{\sqrt{u^2}}
\]
If we choose proper-time parametrization:
\[
u^2 = 1
\]
then:
\[
p_\mu = -m u_\mu
\]
4. Is this a problem?
No — because raising the index gives:
\[
p^\mu = \eta^{\mu\nu} p_\nu = -m u^\mu
\]
But what we physically call momentum is defined as:
\[
p^\mu_{\text{phys}} = m u^\mu
\]
So again, the minus sign is purely conventional.

It corresponds to choosing:
- future-directed vs past-directed parametrization,
- or reversing the orientation of the worldline.
5. What about the Hamiltonian?
Compute:
\[
H = p_\mu u^\mu - L
\]
\[
= (-m u_\mu u^\mu) - (-m \sqrt{u^2})
= -m u^2 + m \sqrt{u^2}
\]
For \(u^2 = 1\):
\[
H = -m + m = 0
\]
Again: Hamiltonian vanishes due to reparametrization invariance.
6. Mass-shell constraint
From the definition:
\[
p_\mu p^\mu = m^2
\]
which is the correct mass-shell condition for signature \((+,-,-,-)\).

Compare with the other convention:
Metric	Constraint
\((-,+,+,+)\)	\(p^2 + m^2 = 0\)
\((+,-,-,-)\)	\(p^2 - m^2 = 0\)


Same physics, different sign convention.
7. Relation to field theory
Scalar field Lagrangian with this signature:
\[
\mathcal{L}
= \frac{1}{2} \partial_\mu \phi \partial^\mu \phi - \frac{1}{2} m^2 \phi^2
\]
Expands to:
\[
\mathcal{L}
= \frac{1}{2} (\partial_t \phi)^2
- \frac{1}{2} (\nabla \phi)^2
- \frac{1}{2} m^2 \phi^2
\]
Same physics as before.
8. Summary table
Feature	Signature (+−−−)	Signature (−+++)
Timelike norm	\(u^2 > 0\)	\(u^2 < 0\)
Action	\(-m\sqrt{u^2}\)	\(-m\sqrt{-u^2}\)
Constraint	\(p^2 = m^2\)	\(p^2 = -m^2\)
Energy	\(p^0 > 0\)	\(p^0 > 0\)
Physics	identical	identical


9. Conceptual takeaway
The metric signature only decides where the minus sign lives — in the metric or in the constraint — but not the physics.

The “wrong sign” only appears when you forget that:
- \(p_\mu\) is a covector,
- raising/lowering indices changes signs,
- reparametrization invariance removes absolute normalization.
10. One-line conclusion
Changing the metric signature reshuffles signs between the Lagrangian, momentum, and constraint, but leaves the physics completely invariant.
1. What “nonrelativistic limit” really means
The nonrelativistic (NR) limit is not simply “small velocity”.

It is a gauge fixing + expansion:
1. Choose a preferred time coordinate
2. Break reparametrization invariance
3. Expand in \(v/c \ll 1\)
4. Drop terms of order \(v^4/c^4\)
This already hides most relativistic structure.
2. Start from the fully covariant action
Take the relativistic point particle:
\[
S = -m \int d\lambda \sqrt{u^2}
\quad\text{with}\quad
u^\mu = \frac{dx^\mu}{d\lambda}
\]
Metric signature: \((+,-,-,-)\)
3. Fix the gauge: choose coordinate time
Choose the gauge:
\[
\lambda = t = x^0
\]
Then:
\[
u^\mu = (1, \dot{\mathbf{x}})
\]
and
\[
u^2 = 1 - \dot{\mathbf{x}}^2
\]
This breaks reparametrization invariance and selects a preferred time.
4. Expand the square root
\[
L = -m \sqrt{1 - \dot{\mathbf{x}}^2}
\]
For small velocities:
\[
\sqrt{1 - v^2}
= 1 - \frac{1}{2}v^2 - \frac{1}{8}v^4 + \cdots
\]
So:
\[
L = -m + \frac{1}{2} m v^2 + \mathcal{O}(v^4)
\]
5. The key moment: discarding the rest energy
The constant term \(-m\) is dropped because it:

• does not affect equations of motion
• does not affect canonical momenta
• does not affect Poisson brackets  

This is where the relativistic structure disappears.

We redefine:
\[
L_{\text{NR}} = \frac{1}{2} m v^2
\]
This is not innocent — it erases information.
6. What disappears when we do this
Relativistic feature	What happens in NR limit
Reparametrization invariance	Fixed by \(t = x^0\)
Hamiltonian constraint	Becomes trivial
Mass-shell condition	Becomes definition of energy
Lorentz symmetry	Reduced to Galilean
Proper time	Identified with coordinate time
Negative metric component	Disappears
Square-root action	Linearized


7. Canonical momenta comparison
Relativistic:
\[
p_i = \frac{\partial L}{\partial \dot{x}^i}
= \frac{m \dot{x}^i}{\sqrt{1 - v^2}}
\]
Nonrelativistic limit:
\[
p_i \approx m \dot{x}^i
\]
Time component:
\[
p_0 = -m \sqrt{1 - v^2} \approx -m + \frac{1}{2}mv^2
\]
The huge constant \(-m\) is simply discarded.
8. Energy emerges from a constraint
Relativistically:
\[
p_\mu p^\mu = m^2
\]
Solve for energy:
\[
p_0 = \sqrt{m^2 + \mathbf{p}^2}
\]
Expand:
\[
p_0 = m + \frac{\mathbf{p}^2}{2m} + \cdots
\]
Subtract the constant \(m\):
\[
E_{\text{NR}} = \frac{\mathbf{p}^2}{2m}
\]
So energy is not fundamental — it is extracted from the constraint after gauge fixing.
9. Why the sign problem disappears
In the NR limit:
- There is no distinction between covariant and contravariant time components
- There is no Lorentzian signature to worry about
- There is no square root
- There is no constraint surface
Everything becomes Euclidean and positive definite.

That’s why the “wrong sign” problem simply vanishes.
10. Geometric interpretation
Relativistic mechanics:
Motion = geometry of worldlines in spacetime

Nonrelativistic mechanics:
Motion = evolution in absolute time

By choosing a preferred time, you collapse spacetime geometry into dynamics.
11. One-sentence summary
The nonrelativistic limit hides the sign issues because it gauge-fixes time, eliminates the Hamiltonian constraint, and discards the rest energy that carries all information about Lorentzian geometry.

12. Final intuition
Relativity:
“Time is part of geometry.”

Nonrelativistic mechanics:
“Time is a background parameter.”

The price of that simplification is that all the beautiful sign structure disappears2. Quadratic (einbein) action
The quadratic relativistic action is
\[
S = \frac{1}{2} \int d\lambda 
\left(
e^{-1} \dot x^\mu \dot x_\mu - e m^2
\right)
\]
This action is:
- reparametrization invariant
- Lorentz invariant
- equivalent to the square-root action
3. Equations of motion
(a) Variation w.r.t. \(x^\mu\)
\[
\delta S = \int d\lambda\, e^{-1} \dot x_\mu \delta \dot x^\mu
\]
Integrating by parts:
\[
\frac{d}{d\lambda} (e^{-1} \dot x_\mu) = 0
\]
Define canonical momentum:
\[
p_\mu := \frac{\partial L}{\partial \dot x^\mu} = e^{-1} \dot x_\mu
\]
(b) Variation w.r.t. the einbein \(e\)
\[
\frac{\delta S}{\delta e} = 0
\quad\Rightarrow\quad
- e^{-2} \dot x^2 - m^2 = 0
\]
So:
\[
\dot x^2 + e^2 m^2 = 0
\]
This is the mass-shell constraint.
4. Solving the constraint
Since \(\dot x^2 > 0\) (timelike):
\[
e = \frac{\sqrt{\dot x^2}}{m}
\]
Plug back into the action:
\[
S = \frac{1}{2} \int d\lambda
\left(
\frac{\dot x^2}{\sqrt{\dot x^2}/m}
- m^2 \frac{\sqrt{\dot x^2}}{m}
\right)
= - m \int d\lambda \sqrt{\dot x^2}
\]
which is exactly the usual square-root action (with this signature).
5. Hamiltonian formulation
Canonical Hamiltonian:
\[
H = p_\mu \dot x^\mu - L
\]
Substitute:
\[
p_\mu = e^{-1} \dot x_\mu
\]
\[
H = e \left( \frac{1}{2} p^2 + \frac{1}{2} m^2 \right)
\]
where
\[
p^2 = \eta^{\mu\nu} p_\mu p_\nu
\]
6. Constraint structure
Because \(e\) has no momentum:
\[
\pi_e = 0
\]
and its equation of motion gives:
\[
\mathcal{H} \equiv \frac{1}{2}(p^2 + m^2) \approx 0
\]
This is the mass-shell constraint.

The Hamiltonian is pure constraint:
\[
H = e \mathcal{H}
\]
7. Gauge symmetry (reparametrization invariance)
Under:
\[
\lambda \to \lambda'(\lambda)
\]
the einbein transforms as:
\[
e(\lambda) \to e'(\lambda') = \frac{d\lambda}{d\lambda'} e(\lambda)
\]
This guarantees invariance of the action.
8. Gauge fixing
Proper time gauge:
\[
e = \frac{1}{m}
\Rightarrow
\dot x^2 = 1
\]
Then:
\[
p^\mu = m \dot x^\mu
\]
Coordinate time gauge:
\[
x^0 = \lambda
\]
Then:
\[
p^0 = \sqrt{m^2 + \mathbf{p}^2}
\]
and:
\[
H = p^0
\]
9. Nonrelativistic limit
Expand:
\[
p^0 = \sqrt{m^2 + \mathbf{p}^2}
= m + \frac{\mathbf{p}^2}{2m} - \frac{\mathbf{p}^4}{8m^3} + \cdots
\]
Drop constant \(m\):
\[
H_{\text{NR}} = \frac{\mathbf{p}^2}{2m}
\]
All relativistic structure disappears.
10. Summary table
Concept	Expression
Action	\(\frac12 \int (e^{-1}\dot x^2 - e m^2) d\lambda\)
Constraint	\(p^2 + m^2 = 0\)
Gauge symmetry	reparametrization
Physical Hamiltonian	emerges after gauge fixing
Metric	\((+,-,-,-)\)


11. Final takeaway
With signature \((+,-,-,-)\), the quadratic relativistic Lagrangian cleanly shows that dynamics is entirely encoded in a constraint; time evolution appears only after gauge fixing, and all sign issues are bookkeeping artifacts of Lorentzian geometry..

